\documentclass[11pt]{article}
\usepackage[a4paper,margin=25mm]{geometry}
\usepackage{amsmath,amssymb,hyperref}
\title{A missing symmetry hypothesis in Hutchinson's variance law}
\author{ziangni-sys}
\date{Conjecture 00000008782}
\begin{document}
\raggedright
\maketitle
\noindent\textbf{Result.} The displayed variance formula is false for general real matrices. A real skew-symmetric $2\times2$ matrix gives variance zero, whereas the proposed expression equals four.

\paragraph{Scope of the original statement.}
Both the English and Chinese versions assert that Hutchinson's variance is
\[
2\left(\|A\|_F^2-\sum_i a_{ii}^2\right).
\]
Neither version states that $A$ must be symmetric.
We use the ordinary one-sample Hutchinson estimator $H_A(z)=z^{\mathsf T}Az$, where the coordinates of $z$ are independent Rademacher variables.
This counterexample concerns the variance conjunct of the unqualified statement. It does not contradict the familiar formula for symmetric real matrices and does not settle the separate Hutch++ or probing clauses.

\paragraph{The actual probability space and matrix.}
Let $\Omega=\{0,1\}^2$ with all subsets measurable and every outcome of probability $1/4$. Put $s(0)=-1$, $s(1)=1$ and $z(\omega)=(s(\omega_1),s(\omega_2))$. Each coordinate takes $-1$ and $1$ with probability $1/2$. For all subsets $S,T\subseteq\mathbb R$,
\[
\mathbb P(z_1\in S,z_2\in T)
=\frac{\#(S\cap\{-1,1\})\#(T\cap\{-1,1\})}{4}
=\mathbb P(z_1\in S)\mathbb P(z_2\in T),
\]
so these are genuinely independent Rademacher coordinates. Take
\[
A=\begin{pmatrix}0&1\\-1&0\end{pmatrix}.
\]
For every real vector $v$, not just the four sampled vectors,
\[
v^{\mathsf T}Av=v_1v_2-v_2v_1=0.
\]
Consequently $H_A$ is identically zero, $\mathbb E H_A=0=\operatorname{tr}A$, and $\operatorname{Var}(H_A)=0$. There is no integrability issue: the estimator is a constant on a finite probability space.

\newpage
\paragraph{The contradicted numerical value.}
The Frobenius norm is defined by $\|A\|_F=\sqrt{\sum_{i,j}|a_{ij}|^2}$. In this real example,
\[
\|A\|_F^2=0^2+1^2+(-1)^2+0^2=2,
\qquad \sum_i a_{ii}^2=0.
\]
Thus the conjectured right side is $2(2-0)=4$, different from the actual variance zero. The original conjunction is therefore false as written.

\paragraph{Why symmetry matters.}
For a general real matrix and independent Rademacher coordinates, expand the estimator as
\[
H_A(z)=\operatorname{tr}A+\sum_{i<j}(a_{ij}+a_{ji})z_i z_j.
\]
The products in this sum have mean zero. Distinct unordered pairs give orthogonal products in $L^2$: some coordinate appears to an odd power, so independence and zero coordinate means make their cross expectation vanish. Therefore
\[
\operatorname{Var}(H_A)=\sum_{i<j}(a_{ij}+a_{ji})^2.
\]
If $A$ is symmetric, this is $4\sum_{i<j}a_{ij}^2=2(\|A\|_F^2-\sum_i a_{ii}^2)$. If $A$ is skew-symmetric, every off-diagonal sum is zero. This explains precisely the missing assumption. This explanatory general formula is not needed as an assumption or as an imported theorem in the formal counterexample.

\paragraph{Formal verification.}
The accompanying Lean 4.19.0 project uses pinned Mathlib commit \texttt{c44e0c8ee63ca166450922a373c7409c5d26b00b}. It constructs the uniform four-outcome probability mass function and its actual probability measure, proves \texttt{IndepFun} of the two real coordinates for all measurable events, and computes their Rademacher distributions. It defines the actual matrix and the finite-sum quadratic form, proves that this form vanishes for every real vector, and computes the Bochner expectation, actual \texttt{ProbabilityTheory.variance}, matrix trace, Frobenius norm, and diagonal square sum. The final theorem \texttt{variance\_formula\_false} proves the displayed inequality of the two expressions. Dependency audits use only Lean's standard logical axioms; no custom axioms, admitted facts or external decision procedures are used.
\end{document}
